\title{Why Rust is the Future of Embedded Development}
\date{2025-12-08}
\excerpt{Exploring how Rust's ownership model and zero-cost abstractions make it an excellent choice for safety-critical embedded systems.}
\tags{Rust, Embedded, Safety}
\xpost{https://x.com/YOUR_HANDLE/status/PLACEHOLDER_TWEET_ID_3}

\begin{document}

\section{The Problem with C in Embedded}

C has been the lingua franca of embedded development for decades, but it comes with well-known pitfalls: buffer overflows, null pointer dereferences, and data races. In safety-critical systems, these bugs can be catastrophic.

\section{Enter Rust}

Rust's ownership model prevents entire classes of bugs at compile time:

\begin{itemize}
  \item \textbf{No null pointers} — Option types instead
  \item \textbf{No data races} — ownership and borrowing rules
  \item \textbf{No buffer overflows} — bounds checking
  \item \textbf{Zero-cost abstractions} — no runtime overhead
\end{itemize}

\begin{lstlisting}[language=rust]
fn read_sensor(peripheral: &mut SensorPeripheral) -> Result<f32, SensorError> {
    let raw = peripheral.read_register(TEMP_REG)?;
    Ok(raw as f32 * 0.0625)
}
\end{lstlisting}

\section{The Ecosystem}

The embedded Rust ecosystem has matured significantly with projects like \texttt{embassy}, \texttt{probe-rs}, and hardware abstraction layers for most popular MCU families.

\section{Conclusion}

While C isn't going anywhere soon, Rust offers a compelling path forward for new embedded projects, especially those with safety or reliability requirements.

\end{document}

